% =====================================================================
% Rubric Deliverable T1 — AI Technical Depth & Non-Triviality
% MRI-ReportGenerator (EECE503N). Self-contained; compile with pdfLaTeX.
% =====================================================================
\documentclass[11pt]{article}
\usepackage[margin=1in]{geometry}
\usepackage{booktabs}
\usepackage{tabularx}
\usepackage{array}
\usepackage{enumitem}
\usepackage{xcolor}
\usepackage{verbatim}
\usepackage{textcomp}
\usepackage[hidelinks]{hyperref}

\newcolumntype{Y}{>{\raggedright\arraybackslash}X}
\setlist{nosep, leftmargin=1.4em}

\title{\textbf{AI Technical Depth \& Non-Triviality}\\[2pt]
\large Rubric Deliverable T1 --- MRI-ReportGenerator (Cervical Spine MRI Analysis)}
\author{EECE503N Final Project}
\date{2026-06-08}

\begin{document}
\maketitle

\noindent\textbf{Scope.} This document argues, with evidence, that the MRI-ReportGenerator pipeline
carries genuine technical and AI depth and is \emph{not} a thin wrapper over pretrained models. It is
written for a skeptical technical reader. Every clinical claim is cited; every engineering claim is
backed by committed code, validation results, and the development journal. Companion evidence:
\texttt{DEVELOPMENT\_JOURNEY.md} (J1--J23), \texttt{docs/validation/group-status-2026-06-08.md}
(per-group verdict), \texttt{docs/validation/results-full-2026-06-08.md} (numbers + figures), and the
companion \texttt{paper/}.

\medskip
\noindent\textbf{Medical framing (hard rule).} The system produces measurements and \emph{findings
flagged for physician review}, never diagnoses. Nothing here claims otherwise.

% ---------------------------------------------------------------------
\section{Thesis}
The depth of this system is not in training a new neural network --- and we deliberately do \textbf{not}
claim to. It is in three places a ``wrapper'' does not reach:

\begin{enumerate}
  \item \textbf{Multi-model segmentation composition} --- orchestrating three independent,
  differently-trained segmentation engines (TotalSpineSeg, Spinal Cord Toolbox, SPINEPS) and deciding,
  per measurement, which one to trust on a region (cervical spine) where none is individually validated
  for the downstream measurement.
  \item \textbf{Clinically-meaningful measurement algorithms} --- converting voxel masks into vertebral,
  disc, canal, cord and alignment measurements with methods that required real geometric engineering
  (endplate-line fitting, anatomical canal-cut body isolation, PCA tilt-correction, endplate-voxel
  Cobb). The na\"ive version of each \emph{demonstrably fails}; we have the before/after numbers.
  \item \textbf{Validation methodology under a hard data constraint} --- no public dataset pairs
  cervical MRI with per-case expert measurements, so per-case sensitivity/specificity is impossible. We
  designed a \emph{threshold-crossing / distribution-separation} methodology anchored on a healthy
  cohort, with explicit confound control. Designing a defensible evaluation when the obvious one is
  unavailable is itself the contribution.
\end{enumerate}

The single strongest piece of evidence that this is non-trivial is the \textbf{bug ledger}
(Section~\ref{sec:ledger}): places where the obvious implementation produced \emph{wrong} clinical
numbers, which only rigorous validation surfaced and principled algorithm changes fixed. A trivial
wrapper has no such ledger.

% ---------------------------------------------------------------------
\section{The honest framing: where the AI is, and where it is not}
We are explicit about this because the rubric rewards honesty and because over-claiming in medical AI is
a failure mode in itself.

\paragraph{What we do \emph{not} do.} We do not train a segmentation model from scratch and do not claim
a novel deep-learning architecture. The segmenters are pretrained and \textbf{frozen}.

\paragraph{Why that is the correct engineering decision, not a shortcut.}
\begin{itemize}
  \item \textbf{Reproducibility and no overfitting.} Because nothing is trained on our evaluation
  cohort, there is no train/test leakage and no risk of fitting our metrics to the symptomatic set. The
  unhealthy cohort is a \emph{demonstration} set, not a training set --- a deliberate medical-AI safety
  property.
  \item \textbf{Standing on validated tools.} TotalSpineSeg and Spinal Cord Toolbox are published,
  peer-reviewed, open-source segmentation stacks. Re-implementing inferior versions would be worse
  engineering, not better.
  \item \textbf{The value lives in the layer they do not provide.} None of these tools outputs a
  cervical radiology report, threshold-based interpretation, or validated cervical \emph{measurements}.
  That layer --- the clinically actionable output --- is entirely ours, and is where the difficulty is.
\end{itemize}

The correct mental model is \textbf{not} ``wrapper around a model'' but ``a
measurement-and-interpretation system that \emph{uses} segmentation as one input, the way a radiology
workstation uses a DICOM loader.''

% ---------------------------------------------------------------------
\section{Pipeline overview}
\begin{verbatim}
Input (DICOM / NIfTI, sagittal T2)
   |
   v
Segmentation IEP   -- TotalSpineSeg  -> vertebrae, discs, spinal canal (labelmap)
                   -- Spinal Cord Toolbox -> cord & canal cross-sections (per-slice)
                   -- SPINEPS         -> vertebra instances + endplate sheets
   |
   v
Measurements IEP   -- G1 vertebral morphometry (Ha/Hp, AP width, tilt)
                   -- G2 disc (height index, AP spread, signal)
                   -- G3 canal / SAC / cord diameters
                   -- G4 cervical alignment (C3-C7 Cobb)
                   -- G5 fracture screen + myelomalacia screen
   |
   v
Interpretation     -- G6 threshold catalog -> status per finding (cited thresholds)
   |
   v
Reporting IEP      -- structured, radiologist-style report ("flagged for physician review")
\end{verbatim}

The Enterprise Execution Plane (EEP) orchestrates the two deployed Internal Execution Planes
(measurements, reporting); see the infra documentation for the service architecture. This document
concerns the \emph{measurement-science depth} inside the measurements IEP and the segmentation
composition that feeds it.

% ---------------------------------------------------------------------
\section{Depth dimension 1 --- Multi-model segmentation composition}
A wrapper calls one model and trusts it. We call \textbf{three}, because no single one is sufficient for
cervical measurement, and the composition is non-trivial.

\begin{table}[h]
\centering\small
\begin{tabularx}{\textwidth}{@{}l l Y Y@{}}
\toprule
\textbf{Engine} & \textbf{License} & \textbf{What it gives us} & \textbf{Why we need it} \\
\midrule
TotalSpineSeg (TSS) & LGPLv3 & Whole-vertebra, disc and canal labels on an isotropic grid & Backbone
labelmap for G1/G2 and the canal for body isolation \\
Spinal Cord Toolbox (SCT) & LGPLv3 & Per-slice cord and canal cross-sectional diameters & G3 ---
cord/canal/SAC; TSS's canal label is coarse for AP diameters \\
SPINEPS & Apache-2.0 & Per-vertebra instances \emph{plus} thin endplate voxel sheets & G4 --- the
endplate-voxel Cobb that TSS corners cannot do reliably at C7 \\
\bottomrule
\end{tabularx}
\end{table}

\noindent The non-trivial engineering here:
\begin{itemize}
  \item \textbf{Per-measurement model selection.} The C3--C7 Cobb angle is \emph{more accurate} from
  SPINEPS endplate voxels than from the TSS labelmap, because the cervico-thoracic junction (C7) is
  frequently obscured and corner-based angles there are unstable. So \texttt{c3c7\_cobb\_angle}
  \textbf{prefers} the SPINEPS method when a seg-vert mask is present and \textbf{falls back} to the TSS
  canal-cut method otherwise. This selection logic measurably rescues cases the fallback cannot measure
  at all (J22).
  \item \textbf{Reconciling incompatible runtimes.} TSS/nnU-Net and SPINEPS pin \textbf{mutually
  incompatible NumPy versions} (SPINEPS requires \texttt{numpy==2.0.2}, which breaks the TSS/nnU-Net
  ABI). They cannot share a process; the pipeline runs them in separate environments and re-aligns their
  outputs onto a common grid.
  \item \textbf{Grid and orientation reconciliation.} The three engines emit masks in different
  orientations/resolutions; \texttt{load\_context} standardizes to canonical RAS, enforces isotropy, and
  index-aligns the grayscale and the SPINEPS instance mask to the labelmap so measurements are computed
  in one coherent frame.
\end{itemize}

None of this is visible to a user, and none of it is provided by any single model. It is the connective
tissue that makes three research tools into one measurement system.

% ---------------------------------------------------------------------
\section{Depth dimension 2 --- Measurement algorithms (the core contribution)}
For each measurement we did \emph{not} take the obvious pixel-counting approach, because the obvious
approach gives clinically wrong numbers on real cervical anatomy. Four representative examples, each with
the evidence that the na\"ive method fails:

\subsection{Endplate-LINE fitting instead of corner extrema (G1 heights, G4 Cobb)}
\textbf{Na\"ive:} read heights from the four corner pixels. \textbf{Why it fails:} cervical endplates
are concave and sloped, not flat~\cite{chen2013}, so a single corner pixel lands wrong.
\textbf{Our method:} fit straight lines (Theil--Sen, robust to osteophyte tails) to the full superior
and inferior endplate boundaries after PCA orientation, and read wall heights from the lines. The
literature supports this directly: a fitted endplate line reaches ICC $\approx 0.97$ versus $\approx
0.75$ for four-corner extraction~\cite{wang2023}.

\textbf{Evidence --- measured, not asserted:} the corner method reads healthy cervical Ha/Hp $\approx
\mathbf{1.08}$ (anterior taller than posterior --- physiologically backwards). The endplate-line method
reads $\mathbf{0.93}$, matching the physiological slight anterior wedge. Switching the service to the
line fit moved healthy Ha/Hp $1.08 \to 0.93$ and preserved the correct healthy $\geq$ unhealthy ordering
(J20, J22).

\subsection{Anatomical canal-cut body isolation instead of morphological erosion (G1)}
\textbf{Na\"ive:} isolate the body by erosion / connected components. \textbf{Why it fails:} on real
cervical anatomy the body and posterior arch are one connected blob on the mid-sagittal slice; erosion
leaves a wedge-shaped body and a systematic false anterior wedge. \textbf{Our method:} cut the mask at
the spinal canal's anterior face --- the body is, by definition, everything anterior to the canal ---
using TSS's own canal label. This replaced a heuristic that produced a systematic $\sim 0.86$
height-ratio artifact.

\subsection{SPINEPS endplate-voxel Cobb instead of labelmap corners (G4)}
\textbf{Na\"ive:} Cobb from the C3 and C7 corner points of the labelmap. \textbf{Why it fails:} C7 is
frequently obscured at the cervico-thoracic junction; corner-based C6--C7 angles have a standard
deviation of $\sim 18.5^\circ$, uselessly noisy. \textbf{Our method:} fit the line directly to SPINEPS'
own inferior-endplate voxel sheets, dropping C6--C7 SD to $\mathbf{5.9^\circ}$ (J11, J12). On the
healthy cohort this reads $+15.2^\circ$ lordosis, matching the external reference (F1000 cervical
cohort, $15.4^\circ$;~\cite{f1000}).

\subsection{Sub-voxel boundary refinement and PCA tilt-correction}
\begin{itemize}
  \item \textbf{Sub-voxel refinement:} the in-plane mask boundary is upsampled and re-thresholded at the
  0.5 level (a marching-squares-style half-level crossing), placing edges between voxel centres rather
  than on them, reducing the $\pm 1$-voxel quantization on every height/width/slip.
  \item \textbf{PCA tilt-correction:} every vertebra is measured on its \emph{own} principal axes,
  because healthy cervical bodies sit $\sim 28^\circ$ from absolute vertical (the lordotic curve).
  Measuring on image axes would conflate normal lordosis with deformity --- exactly the bug the old
  \texttt{tilt} flag had (Section~\ref{sec:ledger}).
\end{itemize}

Each is a deliberate algorithmic choice with a measured payoff. That is the definition of technical depth
in a measurement pipeline.

% ---------------------------------------------------------------------
\section{Depth dimension 3 --- Validation methodology under a hard data constraint}
The hardest intellectual problem was not building the measurements --- it was figuring out how to
\textbf{validate} them when the ideal data does not exist.

\paragraph{The constraint.} Extensive, repeated dataset searches established that \textbf{no public
dataset pairs cervical MRI with per-case expert measurements.} Per-case sensitivity/specificity against a
radiologist gold standard is therefore impossible with public data.

\paragraph{The methodology we designed instead.} Our measurements are \emph{disease-agnostic
geometry/signal detectors anchored on healthy norms} --- a canal of 8\,mm reads ``stenosis'' regardless
of cause. This yields a validation strategy that needs no per-case labels:
\begin{enumerate}
  \item \textbf{Healthy-anchoring.} Validate that healthy anatomy reads \emph{inside} the normal range
  (specificity), using a healthy cohort (Spine-Generic).
  \item \textbf{Threshold-crossing / distribution separation.} Validate that symptomatic anatomy
  \emph{crosses} into the abnormal range (the symptomatic MMCSD cohort), via distribution-separation
  statistics (Mann--Whitney), not per-case accuracy.
  \item \textbf{The scanner-immunity insight.} We discovered and exploited a key distinction:
  \emph{physical-dimension} metrics (mm --- canal, SAC, cord) are immune to scanner differences and
  validate cross-dataset, while \emph{intensity/ratio} metrics (disc signal, height ratios) are
  acquisition-sensitive and must be validated \emph{within} a single dataset. This is why G3 (mm)
  validates cleanly cross-dataset and the disc signal does not --- a non-obvious property that shaped the
  entire evaluation design (J16, J23).
  \item \textbf{Confound control.} When the within-dataset disc test showed a large effect for disc AP
  width (AUC 0.79), we did \emph{not} report it --- lesion discs cluster at wider mid-cervical levels, so
  we re-tested level-stratified and the honest figure was AUC 0.61. Reporting the unstratified 0.79 would
  have been a confounded over-claim (J23).
\end{enumerate}

Designing, justifying, and executing a sound evaluation when the textbook one is unavailable is a
graduate-level methods contribution, not a wrapper feature.

% ---------------------------------------------------------------------
\section{Proof of non-triviality: the bug ledger}
\label{sec:ledger}
The most direct evidence that this pipeline required real depth is the set of clinically-wrong outputs
the \emph{na\"ive} implementation produced and that only rigorous validation caught. If the task were
trivial, these would not exist.

\begin{table}[h]
\centering\small
\begin{tabularx}{\textwidth}{@{}Y Y Y l@{}}
\toprule
\textbf{Na\"ive behavior (wrong)} & \textbf{How it was caught} & \textbf{Principled fix} &
\textbf{Evidence} \\
\midrule
Vertebral Ha/Hp read \emph{backwards} (anterior taller, 1.08) & Healthy cohort read non-physiological &
Endplate-line fit replaces corner extrema & $1.08\!\to\!0.93$ \\
\texttt{tilt} flag fired on \emph{88\% of healthy} vertebrae & Healthy false-flag rate measured &
Recalibrated cut $20^\circ\!\to\!45^\circ$ (lordosis is normal) & $88\%\!\to\!0\%$ \\
Disc posterior bulge \emph{backwards} (healthy worse, 60\% flagged) & Healthy over-flagged vs
symptomatic & Reference chord from endplate-line corners & $60\%\!\to\!8\%$ \\
Disc AP-width ``AUC 0.79'' discrimination & Suspected level confound & Level-stratified re-test & Honest
AUC 0.61 \\
Disc signal grade looked usable & Within-dataset lesion test & Shown non-discriminating (AUC 0.50) &
Reported as a negative \\
Latent: \texttt{seg == "C4"} (string vs int label) matched nothing & Code review during bulge fix &
Resolve names via \texttt{VERT\_LABELS} & J22 \\
\bottomrule
\end{tabularx}
\end{table}

\noindent Each row is a place where ``just measure it'' gives the wrong clinical answer, and a place
where we have a committed, tested fix. The full narrative is in \texttt{DEVELOPMENT\_JOURNEY.md}.

% ---------------------------------------------------------------------
\section{Validated outcomes (the depth pays off)}
We are equally explicit about what is strongly validated and what is partial --- honesty is a graded and
a clinical requirement.

\begin{table}[h]
\centering\small
\begin{tabularx}{\textwidth}{@{}l Y l@{}}
\toprule
\textbf{Group} & \textbf{Result} & \textbf{Status} \\
\midrule
G3 canal / SAC / cord (SCT) & canal AP $11.7\!\to\!8.6$\,mm \textbf{p=0.0001}; SAC $4.7\!\to\!2.3$\,mm
\textbf{p=0.0001}; cord $6.3\!\to\!5.5$\,mm p=0.009 & \textbf{Validated (strong)} \\
G4 C3--C7 Cobb (SPINEPS) & Correct lordosis $+15.2^\circ$ ($=$ literature $15.4^\circ$); discrimination
underpowered ($d=0.91$, awaiting larger $n$) & Method-validated \\
G1 Ha/Hp compression screen & 0\% false-flag on healthy; correctly null on (non-compressive)
spondylosis, confirmed $n=49$ & Validated as a screen \\
G2 disc & Signal/bulge documented negatives; disc/VB AP ratio discriminates (AUC 0.62, level-controlled)
& Partial --- geometric spread \\
G5 myelomalacia / fracture screens & $\sim$91\% healthy specificity (myelomalacia); compression screen
as G1 & Validated as screens \\
\bottomrule
\end{tabularx}
\end{table}

% ---------------------------------------------------------------------
\section{Why this composition is the right AI choice for clinical deployment}
The pretrained-and-compose design is the defensible architecture for a clinical \emph{application} (the
project's positioning):
\begin{itemize}
  \item \textbf{Auditability.} Every output traces to a measurement and a \emph{cited} threshold (G6
  catalog), not to an opaque end-to-end model. A radiologist can check each number.
  \item \textbf{No overfitting by construction.} Frozen segmenters + cited thresholds mean the system
  cannot have memorized our evaluation cohort.
  \item \textbf{Graceful degradation.} Per-component error contracts mean one failed measurement does
  not sink the report; a missing SPINEPS mask degrades the Cobb method, it does not crash the pipeline.
  \item \textbf{The differentiator is the clinical layer.} Measurement validity, cited interpretation,
  and the report --- the parts that decide whether a radiologist would trust the output --- are exactly
  the parts we built and validated.
\end{itemize}

% ---------------------------------------------------------------------
\section{Summary for the grader}
\begin{itemize}
  \item The system is \textbf{not} a wrapper: it composes three segmentation engines, implements
  non-trivial geometric measurement algorithms, and validates them under a genuine data constraint.
  \item The depth is \textbf{demonstrated, not asserted} --- every algorithmic choice has a measured
  before/after, and the bug ledger shows the na\"ive approach producing clinically wrong numbers we
  caught and fixed.
  \item We are honest about the boundary: pretrained segmenters (correct medical-AI choice), validated
  measurements (our contribution), and an evaluation methodology designed for a domain where the ideal
  dataset does not exist.
\end{itemize}

% ---------------------------------------------------------------------
\begin{thebibliography}{9}
\bibitem{genant1993} Genant HK, et al. Vertebral fracture assessment using a semiquantitative technique.
\emph{J Bone Miner Res.} 1993. PMID:~8237484.
\bibitem{wang2023} Wang, et al. Endplate line-fit morphometry (ICC 0.97 vs 0.75 four-corner). 2023.
PMC10685593.
\bibitem{chen2013} Chen, et al. Cervical endplate concavity/slope morphology. 2013.
\bibitem{miyazaki2008} Miyazaki M, et al. Cervical disc degeneration grading (T2). 2008. PMID:~18525490.
\bibitem{nakashima} Nakashima H, et al. Disc bulging on MRI in an asymptomatic population ($n=1211$).
PMID:~25584950.
\bibitem{nell2019} Nell C, et al. Cervical canal / SAC normative values on MRI. \emph{PLoS One.} 2019.
doi:10.1371/journal.pone.0222682.
\bibitem{f1000} F1000 cervical MRI + Cobb cohort ($n=77$) --- external alignment reference.
\bibitem{zhang2025} Zhang, et al. Automated cervical Cobb (sagittal-T2 C2--C7), ICC 0.94 / MAE
$2.44^\circ$. 2025.
\end{thebibliography}

\paragraph{Cohorts.} Spine-Generic (healthy multi-vendor cervical T2) --- healthy anchor cohort; MMCSD
(Synapse syn63903115) --- symptomatic CSM/CSR cohort, used as the demonstration set. Threshold
provenance (including in-code values found uncited and corrected): see the G6 threshold catalog and the
verified-research memos.

\end{document}
